\begin{abstract}
\textbf{Objective:} The apparent advantage of one EEG covariance
regularisation selector can depend on what is actually measured
in held-source validation. We distinguish covariance fitting
after validation-person class re-centring from proper
source-mean plug-in Gaussian predictive log scoring.
\textbf{Methods:} We audit previously frozen zero-target
cross-subject EEG decoding studies, derive the exact
nonnegative class-mean-shift term omitted by the historical
covariance score, and calculate exact source-only selection
breakpoints along a parametric continuation between
the two objectives. Corrected predictive controls
were run retrospectively on two independently archived
public EEG cohorts and checked against their immutable
target AUC paths.
\textbf{Results:} The original prospective EEGMMIDB
participant block ($N=20$) gave a $+0.0561$ mean
AUC contrast between source discrimination and
historical class-centred covariance scoring
(Holm $p=0.00597$), with five materially harmed
individuals. Two originally frozen BNCI tests were
individually non-significant and a prospective
filter-bank gain hypothesis was contradicted.
In retrospective Shin2017A ($N=29$), discrimination,
corrected predictive and historical covariance
selectors achieved mean AUCs of 0.785900,
0.786207 and 0.763678; in retrospective
BNCI2014\_002 ($N=14$), the respective
values were 0.778705, 0.771886 and 0.761116.
The median exact source-only weight at which
the historical selection first changes was
0.0660 and 0.1944 in the two cohorts.
\textbf{Conclusion:} Held-source participant class
re-centring suppresses a mathematically
necessary between-person mean-shift term and
can change the estimated best regularisation.
The two retrospective corrected controls do not
demonstrate universal superiority or equivalence.
Estimand-specific reporting and independent
pre-data transfer evaluation remain necessary.
\end{abstract}
